% 第六章 语音交互系统设计

\section{语音交互系统设计}

\subsection{系统架构}
\par 语音交互子系统的目标是在不改变硬件链路的前提下，为智能家居平台提供自然语言入口，使用户能够以口语化表达完成"查询环境数据、获取统计摘要、查询天气、记录偏好与重要信息"等任务。系统整体采用模块化流水线架构，由语音识别（ASR）、大语言模型（LLM）、语音合成（TTS）、工具模块与音频I/O模块共同组成，并在运行时通过线程与队列完成解耦调度。
\par 在数据流上，系统将麦克风采集到的PCM音频流送入ASR模块进行实时转写，得到的最终文本进入LLM模块生成回复。LLM模块支持函数式工具调用，可在对话过程中根据用户意图自动选择工具并生成结构化参数，工具模块执行具体的外部请求或本地操作后将结果返回给LLM，LLM再综合结果生成最终自然语言回复。回复文本经TTS模块流式合成为音频片段，音频I/O模块负责播放输出并在播放期间暂停录音以避免回声干扰。该架构既保证交互的自然性，又通过工具执行边界保证系统操作的可控性。
\par 与服务器系统的集成主要通过HTTP接口完成。工具模块内置对智能家居平台REST接口的访问能力，可获取最新环境数据、统计摘要、历史列表与时间范围查询等结果，从而实现"语音问答—平台数据—语音播报"的闭环。对于涉及敏感操作的能力，本系统优先选择只读查询类工具，以降低误触发风险；若后续需要扩展到设备控制类工具，可在工具定义与参数校验层面增加权限与安全约束。

\begin{figure}[H]
\centering
\fbox{\parbox{0.92\linewidth}{\centering 图6-1占位：语音交互系统整体架构图。\\
要求展示模块化流水线架构，包括音频I/O模块、ASR模块、LLM模块、TTS模块、工具模块，并用箭头标注数据流方向（PCM音频流→ASR→文本→LLM→工具调用→工具执行→LLM→文本→TTS→音频流→播放）。同时标注线程与队列的调度关系，以及工具模块与智能家居平台HTTP接口的交互。}}
\caption{语音交互系统整体架构图（建议配图）}
\label{fig:voice_system_architecture}
\end{figure}

\subsection{语音识别模块（ASR）}
\par ASR模块负责将实时语音流转换为文本。本系统采用基于WebSocket的实时转写模式，音频I/O模块以固定分块大小采集PCM音频并持续推送至ASR服务端。ASR模块通过回调机制分别处理"中间转写文本"与"最终转写文本"，其中最终转写文本作为一次用户输入进入后续LLM处理队列。实时转写模式能够在保证识别准确度的同时降低交互延迟，适合语音助手的连续对话场景。
\par 音频处理侧采用单声道PCM数据流，系统通过分块发送与队列缓冲降低抖动影响。为避免语音播放时麦克风拾音导致的"回声触发"，在LLM开始生成回复并进入TTS播放阶段时，系统会暂停录音并在播放完成后恢复录音，从而提升交互稳定性。

\begin{figure}[H]
\centering
\fbox{\parbox{0.92\linewidth}{\centering 图6-2占位：ASR模块实时转写流程图。\\
要求展示音频采集、分块发送、WebSocket通信、中间转写与最终转写的流程，标注回调机制与文本队列的传递关系。同时展示录音暂停与恢复的时机（TTS播放期间暂停录音）。}}
\caption{ASR模块实时转写流程图（建议配图）}
\label{fig:asr_flow}
\end{figure}

\subsection{大语言模型模块（LLM）}
\par LLM模块负责意图理解、对话管理与工具调用编排。系统在配置层支持对模型名称、是否启用流式输出、系统提示词等参数进行统一管理。LLM模块维护会话历史并在每轮对话中追加用户输入与助手输出，从而保持上下文一致性。为提升连续对话体验，系统支持流式响应并将文本增量输出交给TTS模块进行流式合成，使用户在模型生成过程中即可听到逐段播报，从而降低等待感。
\par 系统提示词设计以"角色约束 + 输出风格约束 + 工具能力声明"为核心。通过在系统提示词中明确助手身份、输出要求与可用工具边界，可降低模型产生不受控行为的概率。系统同时引入记忆模块用于保存用户位置、偏好等长期信息，并在构建系统提示词时动态注入记忆摘要，使模型在后续对话中能够利用已知信息进行更贴近用户的回答。
\par 工具调用机制采用函数式接口描述。工具模块向LLM提供一组结构化的工具定义，包括查询智能家居传感器数据、获取统计摘要、获取历史列表、时间范围查询与导出等接口，以及天气查询与记忆管理等能力。LLM在生成过程中可以输出工具调用请求，系统解析调用名称与JSON参数后执行工具函数，并将执行结果作为工具消息回灌给LLM，最终由LLM生成对用户友好的自然语言回复。该机制将"语言理解"和"外部操作"解耦，有利于提升系统的可靠性与可扩展性。

\begin{figure}[H]
\centering
\fbox{\parbox{0.92\linewidth}{\centering 图6-3占位：LLM工具调用机制流程图。\\
要求展示用户输入→LLM意图理解→工具选择→工具调用请求（函数名+JSON参数）→工具执行→结果返回→LLM综合生成自然语言回复的完整流程。同时展示系统提示词（角色约束+工具能力声明+记忆摘要）的注入位置，以及会话历史的维护机制。}}
\caption{LLM工具调用机制流程图（建议配图）}
\label{fig:llm_tool_calling}
\end{figure}

\subsection{语音合成模块（TTS）}
\par TTS模块负责将LLM输出的文本转换为可播放的语音。本系统采用实时TTS方式，LLM在流式输出文本片段时，TTS模块同步接收并进行流式合成，生成的音频片段进入音频播放队列。为提升播放连续性，系统设置了初始缓冲机制，先累积少量字符再开始合成，避免过短片段导致的频繁启停与听感断裂；随后系统按标点分段进行合成与播放，以提升自然度。
\par 为避免音频队列无限增长导致内存压力，系统在TTS侧与音频输出侧均采用有界队列并实现背压控制。当输出队列使用率过高时，TTS播放线程会短暂等待或丢弃最旧音频块，从而保证系统可持续运行。语速、音色与采样率等参数通过配置统一管理，便于根据设备播放能力与场景需求进行调整。

\begin{figure}[H]
\centering
\fbox{\parbox{0.92\linewidth}{\centering 图6-4占位：TTS流式合成与播放流程图。\\
要求展示LLM流式文本输出→初始缓冲→标点分段→TTS合成→音频片段→有界队列→音频播放的流程。同时标注背压控制机制（队列使用率监控、等待或丢弃策略），以及初始缓冲与分段合成的时机。}}
\caption{TTS流式合成与播放流程图（建议配图）}
\label{fig:tts_flow}
\end{figure}

\subsection{功能工具模块}
\par 工具模块是语音助手"可执行能力"的承载层。环境监测查询工具通过访问智能家居平台API获取最新环境数据、统计摘要、历史数据列表与时间范围查询结果，并将结构化数据格式化为适合语音播报的自然语言描述。为保证交互可读性，工具模块对不同传感器字段配置显示名称与单位映射，并对缺失值进行容错处理，从而避免播报出现不完整或难理解的内容。
\par 天气查询工具作为扩展示例，支持基于用户输入或记忆中保存的位置信息查询实时天气与预报信息。记忆管理工具用于将用户的重要信息持久化保存，并在保存后触发LLM系统提示词刷新，使模型能够在后续对话中利用新增记忆。系统控制接口在当前实现中以数据查询为主，若后续扩展到设备控制类工具，应增加更严格的参数校验、二次确认与权限控制策略，以确保语音控制的安全性与可追溯性。

\begin{figure}[H]
\centering
\fbox{\parbox{0.92\linewidth}{\centering 图6-5占位：工具模块与智能家居平台交互图。\\
要求展示工具模块（环境监测查询、天气查询、记忆管理）与智能家居平台HTTP接口的交互关系，包括API调用、数据格式化、字段映射与容错处理。同时展示记忆管理工具对LLM系统提示词刷新的触发机制。}}
\caption{工具模块与智能家居平台交互图（建议配图）}
\label{fig:tool_module_interaction}
\end{figure}

\par 配图建议如下。图6-6建议为"语音助手获取传感器数据流程与数据通路图"，画法为双流程并行展示：上半部分展示传感器数据采集与上报链路（STM32采集→UART→ESP32网关→HTTP POST→服务器存储），下半部分展示语音助手查询与播报链路（用户语音→ASR→LLM→工具调用→HTTP GET→服务器查询→数据格式化→LLM回复→TTS→音频播放），并用不同颜色区分两条链路。
\begin{figure}[H]
\centering
\fbox{\parbox{0.92\linewidth}{\centering 图6-6占位：语音助手获取传感器数据流程与数据通路图。\\
要求展示两条并行流程：上半部分为传感器数据采集与上报链路（STM32→UART→ESP32→HTTP→服务器存储），下半部分为语音助手查询与播报链路（用户语音→ASR→LLM→工具调用→HTTP GET→服务器查询→格式化→LLM→TTS→播放），并用不同颜色或线型区分。同时标注数据存储与查询的关联关系。}}
\caption{语音助手获取传感器数据流程与数据通路图（建议配图）}
\label{fig:voice_sensor_data_flow}
\end{figure}
